An easy representation about the Turing Machine could be as follows. \vspace{3.5pt}
\begin{center}
    \includegraphics[width=0.5\textwidth]{img/img2.png}
\end{center} \vspace{3.5pt}
where:
\begin{itemize}
    \renewcommand{\labelitemi}{-}
    \item $\{0,1\}^*\rightarrow$ alphabet that provides us a whatever string of any dimension
    \item \textit{Scratchpad} $\rightarrow$ physical part of the Turing Machine
    \item \textit{Instructions} $\rightarrow$ set of symbols used by the Turing Machine
    \item $\{0,1\} \rightarrow$ answer of the Turing Machine as a single character
\end{itemize}
Therefore, given in input a whatever string from the alphabet $\{0,1\}^*$ returns $0$ or $1$, that respectively described when a \textbf{decision problem} is \textbf{decidable}
or \textbf{undecidable}. Informally, if we want to compute a function $f:\{0,1\}^* \rightarrow \{0,1\}$, we need a finite, simple and generic set of instructions.
Instructions of the Turing Machine are divided into:
\begin{enumerate}
    \item Read a bit.
    \item Read a symbol from the  scratchpad.
    \item Write a symbol into the scratchpad.
    \item Declare the computation finished or terminated.
\end{enumerate}
The \textbf{running time} of any Turing Machine on the input $x$ is simply the number of basic instructions made. Generally, we say that the machine \textbf{runs in time}
$T(n)$ if it performs at most $T(n)$ instructions on the input string of length $n$. Notice, every instruction made by the Turing Machine is done for each bit or \textit{cell}
that composed the original string $x$ given in input. \vspace{3.5pt}

Since the set of instructions is finite, the instructions can be encoded as strings, and therefore the machine can be described as a string. 
\begin{definition}
    Given a string $\alpha$, we indicate with $M_{\alpha}$ the Turing Machine encoded by $\alpha$.
\end{definition}
\begin{definition}
    With $U$ we indicate the \textbf{Universal Turing Machine}, which simulates any other Turing Machine given its string representation (\textit{as $\alpha$ in the previous definition}):
    given the pair $(x, \alpha)$ the machine $U$ simulates the behaviour of $M_{\alpha}$ on $x$, where:
    \begin{itemize}
        \renewcommand{\labelitemi}{-}
        \item $\alpha \rightarrow$ encoding of the Turing Machine $M_{\alpha}$
        \item $x \rightarrow$ input string of whenever dimension
    \end{itemize}
\end{definition}
However, there are functions as $f: \{0,1\}^* \rightarrow \{0,1\}$ which are intrinsically uncomputable by Turing Machines. We need a more formal way in order to prove 
their \textbf{uncomputability}.